\documentclass[aps,preprint,showpacs,aps,amsmath,showkeys,nobibnotes,nofootinbib,footnotes,superscriptaddress]{revtex4-2}

\usepackage[utf8]{inputenc}
\usepackage[T2A]{fontenc}
\usepackage[english,russian]{babel}
\usepackage{xcolor}
\usepackage{latexsym,amsmath,amssymb,amsbsy,graphicx}
%\usepackage{svg}
%\DeclareGraphicsExtensions{.eps}

% ====================== new definitions ===========================
\newcommand{\bra}[1]{\mathop{\langle{#1}\left.\hspace{-5 pt}\right|}\nolimits}
\newcommand{\ket}[1]{\mathop{\left|\hspace{-4 pt}\right.{#1}\hspace{-1 pt}\rangle}\nolimits}
% ==============================================================
%%%%%%%%%%%%%%%%%%%%%%%%%%%%%%%%%
\def\andname{\relax и}%   Исправляем "and" на "и" в списке авторов
%%%%%%%%%%%%%%%%%%%%%%%%%%%%%%%%%
\makeatother


\begin{document}

\title{Эффективное моделирование черенковского излучения ШАЛ: модификация CORSIKA для проекта СФЕРА-3}
	
	
	\author{\firstname{Е.~А.}~\surname{Бонвеч}}
	\email{bonvech@yandex.ru}
	\affiliation{Научно-исследовательский институт ядерной физики имени Д.\,В.\,Скобельцына Московского государственного университета имени М.\,В.\,Ломоносова}
    \author{\firstname{М.~Д.}~\surname{Зива}}
	\affiliation{Научно-исследовательский институт ядерной физики имени Д.\,В.\,Скобельцына Московского государственного университета имени М.\,В.\,Ломоносова}
	\affiliation{Факультет вычислительной математики и кибернетики Московского государственного университета имени М.\,В.\,Ломоносова}
	\author{\firstname{В.~И.}~\surname{Галкин}}
	\affiliation{Научно-исследовательский институт ядерной физики имени Д.\,В.\,Скобельцына Московского государственного университета имени М.\,В.\,Ломоносова}
	\affiliation{Физический факультет Московского государственного университета имени М.\,В.\,Ломоносова}
	\author{\firstname{В.~А.}~\surname{Иванов}}
	\affiliation{Научно-исследовательский институт ядерной физики имени Д.\,В.\,Скобельцына Московского государственного университета имени М.\,В.\,Ломоносова}
	\affiliation{Физический факультет Московского государственного университета имени М.\,В.\,Ломоносова}
	\author{\firstname{Т.~А.}~\surname{Колодкин}}
	\affiliation{Научно-исследовательский институт ядерной физики имени Д.\,В.\,Скобельцына Московского государственного университета имени М.\,В.\,Ломоносова}
	\affiliation{Физический факультет Московского государственного университета имени М.\,В.\,Ломоносова}
	\author{\firstname{Н.~О.}~\surname{Овчаренко}}
	\affiliation{Научно-исследовательский институт ядерной физики имени Д.\,В.\,Скобельцына Московского государственного университета имени М.\,В.\,Ломоносова}
	\affiliation{Физический факультет Московского государственного университета имени М.\,В.\,Ломоносова}
	\author{\firstname{Д.~А.}~\surname{Подгрудков}}
	\affiliation{Научно-исследовательский институт ядерной физики имени Д.\,В.\,Скобельцына Московского государственного университета имени М.\,В.\,Ломоносова}
	\affiliation{Физический факультет Московского государственного университета имени М.\,В.\,Ломоносова}
	\author{\firstname{Т.~М.}~\surname{Роганова}}
	\affiliation{Научно-исследовательский институт ядерной физики имени Д.\,В.\,Скобельцына Московского государственного университета имени М.\,В.\,Ломоносова}
	\author{\firstname{О.~В.}~\surname{Черкесова}}
	\affiliation{Научно-исследовательский институт ядерной физики имени Д.\,В.\,Скобельцына Московского государственного университета имени М.\,В.\,Ломоносова}
	\affiliation{Факультет космических исследований Московского государственного университета имени М.~В.\,Ломоносова}
	\author{\firstname{Д.~В.}~\surname{Чернов}}
	\affiliation{Научно-исследовательский институт ядерной физики имени Д.\,В.\,Скобельцына Московского государственного университета имени М.\,В.\,Ломоносова}
	
	
	%\date{\today}
	%\today печатает cегодняшнее число
	
	\begin{abstract}
		...
	\end{abstract}
	
	\maketitle


	УДК 524.1

	Ключевые слова: широкие атмосферные ливни, детектор, черенковский свет, моделирование.


    \section{Введение}

Проблема происхождения и природы первичных космических лучей (ПКЛ) сверхвысоких энергий остается одной из фундаментальных задач современной астрофизики. Несмотря на более чем семидесятилетнюю историю исследований, три классические проблемы — форма энергетического спектра, степень анизотропии потока и химический (массовый) состав частиц — до сих пор не имеют окончательного решения~\cite{KAMPERT2012660}. 

Энергетический спектр ПКЛ охватывает более 12 порядков по величине энергии, что диктует необходимость использования различных методов регистрации. В области энергий ниже $10^{15}$~эВ применяются прямые методы с использованием детекторов на спутниках и высотных аэростатах. Однако при переходе к сверхвысоким энергиям (E > $10^{15}$~эВ) резкое падение потока частиц делает прямые измерения статистически неэффективными, вынуждая исследователей переходить к косвенным методам регистрации широких атмосферного ливня ШАЛ. Особую сложность представляет определение массового состава ПКЛ, так как оно требует прецизионного восстановления параметров каждого зарегистрированного события широкого атмосферного ливня (ШАЛ)~\cite{KAMPERT2012660}.

Метод регистрации ПКЛ сверхвысоких энергий с помощью ШАЛ включает в себя довольно много методик с использованием разнообразных детекторов. В настоящее время мы развиваем черенковскую методику, предложенную А.Е. Чудаковым~\cite{Chudakov1972}, использующую черенковский свет ШАЛ, отраженный от заснеженной поверхности и наблюдаемый с помощью телескопа, установленного на летательном аппарате. Этот метод реализуется в серии экспериментов СФЕРА~\cite{Chernov2015, 2020_SPHERE_JINST, Chernov2024}. Он имеет определенные преимущества перед другими методиками, среди которых дешевизна и мобильность. Особый интерес для нас представляет проблема массового состава ПКЛ, которая является самой сложной из трех упомянутых проблем: она требует точного измерения каждого события ШАЛ для оценки массы первичной частицы~\cite{mass}.

Для оптимизации конструкции детектора и разработки алгоритмов реконструкции событий необходимо проведение масштабного численного моделирования ШАЛ. В качестве основного инструмента используется пакет CORSIKA-7~\cite{CORSIKA-7}, позволяющий детально отслеживать развитие электромагнитных и адронных компонент ливня. Расчеты выполняются на суперкомпьютере «Ломоносов-2»~\cite{SC}. Ключевой проблемой при моделировании событий с энергиями выше 70~ПэВ с генерацией черенковских фотонов является высокое время счета одного события. Оригинальная версия CORSIKA-7 не поддерживает параллельную обработку подкаскадов при генерации черенковского света, что вынуждает использовать одно вычислительное ядро на одно событие. При указанных энергиях время моделирования одного события превышает допустимые лимиты выполнения заданий в пакетной системе суперкомпьютера, что делает невозможным накопление достаточной статистики для анализа массового состава.

В данной работе представлена разработанная нами параллельная версия кода CORSIKA-7, адаптированная для эффективного расчета черенковского излучения ШАЛ на многопроцессорных архитектурах. Мы описываем алгоритм распараллеливания процесса генерации и транспорта черенковских фотонов. Приводятся результаты тестирования производительности новой версии кода на локальном сервере и демонстрируется возможность моделирования событий с энергиями до $10^{17}$~эВ в рамках стандартных временных ограничений, что открывает путь к полномасштабному анализу данных для проекта СФЕРА-3.

\section{Первоначальная адаптация CORSIKA}

С самого начала нашего проекта мы не использовали оригинальный код CORSIKA по крайней мере по двум причинам: а) наша методика не предполагала использование детекторов частиц, и нам не требовались стандартные выходные файлы вторичных частиц CORSIKA; б) мы намеревались создать обширную базу данных искусственных событий ШАЛ с характеристиками черенковского света в соответствии с требованиями проекта для дальнейшего многократного использования. Поэтому нам пришлось модифицировать вывод CORSIKA в соответствии с нашими потребностями. Кроме того, нам также требовалась некоторая дополнительная информация, которая не предоставляется стандартной версией.

После соответствующей модификации кода типичное искусственное событие ШАЛ включает в себя: четырехмерное пространственно-временное распределение черенковского света на уровне снега (озеро Байкал), распределения черенковского света в пространстве, углах и времени прихода на трех высотах над снегом (500, 1000 и 1500~м) в виде гистограмм с разбиением, согласованным с характеристиками детектора. Для более детального изучения процесса генерации черенковского света ШАЛ добавлена возможность сохранять в отдельный файл распределение средних высот излучения черенковских фотонов, пришедших в каждую ячейку многомерного пространственно-углового-временного распределения, описанного выше. 

Эта модификация сделала файлы событий значительно более компактными без потери существенной физической информации о ШАЛ. Одно событие представлено двоичным файлом объемом около 6~ГБ, который сжимается до менее чем 1~ГБ. В настоящее время база данных занимает около 100~ТБ дискового пространства и содержит более $10^5$ событий с энергиями $10^{15}-3\cdot10^{16}$~эВ~\cite{2025NuclPhys}. Для каждой  первичной частицы рассчитывается по 100 событий для каждого набора параметров первичной частицы и моделей расчетов: четыре первичных частицы (протон, ядра гелия, азота, железа), 6 зенитных углов (от 5 до 30 с шагом 5 градусов), четыре модели атмосферы, три модели ядерных взаимодействий. 

Для получения полного набора событий с одной энергией для одной модели ядерного взаимодействия считается 9600 событий. На суперкомпьютере Ломоносов-2 возможно одновременно считать 300 событий на одну задачу. Однако, при моделировании событий с энергиями выше 70~ПэВ время счета одного события на суперкомпьютере Ломоносов-2 превышает допустимые временные лимиты выполнения заданий в двое суток, а при одновременном моделировании даже 10 событий на одном локальном компьютере потребуется несколько лет, чтобы получить полный набор событий для одной модели взаимодействия. Поэтому возникла потребность создать параллельную версию программы, которая бы позволила использовать преимущества вычислений на суперкомпьютере, сократив время вычислений одного события. 


\section{Параллельная версия}

Основная идея нашего подхода к преобразованию модифицированного кода CORSIKA в многопроцессную версию заключалась в достижении двух целей: а) распределить вычислительную нагрузку по трассировке огромного каскада частиц с испусканием черенковских фотонов между несколькими вычислительными ядрами и б) выделить память для многомерных массивов только один раз. Мы разделяем процесс генерации события на два этапа (см. схему работы параллельной программы на рис.~\ref{fig:scheme}). На первом этапе мы выполняем стандартные инициализации CORSIKA и сохраняем эти готовые к использованию данные в двоичный файл \textit{general.bin}.

\begin{figure}[htb]
    \centering
    %\onelinecaptionsfalse % if the caption is multiline
    %\onelinecaptionstrue  % if the caption is one-line
    \includegraphics[width=0.99\textwidth]{figs/corsika_scheme_rus}
    %\includesvg[width=0.7\textwidth]{figs/corsika_scheme_rus}
    \vspace{-2cm}
    \caption{Схема работы параллельной версии CORSIKA}\label{fig:scheme}
\end{figure}

В реализации второго этапа используется технология MPI для создания независимых вычислительных процессов, работающих параллельно. Эта часть программы начинается с чтения файла, после которого основной процесс начинает трассировку первичной частицы через атмосферу, формируя стек вторичных частиц. Здесь мы модифицируем стандартный случайный способ заполнения стека в CORSIKA таким образом, чтобы лидирующая частица (наиболее энергичная частица среди вторичных частиц, рожденных во взаимодействии) записывалась в стек последней и, следовательно, покидала его первой. Такой порядок формирования стека направлен на то, чтобы главный процесс в первую очередь выполнял трассировку лидера. Эта операция относительно быстра и генерирует большой пул вторичных частиц, которые затем могут быть равномерно распределены между вторичными процессами.

Количество взаимодействий лидера, отслеживаемых главным процессом, зависит от массы ядра-снаряда, поскольку взаимодействия ядро-ядро часто оказываются периферийными с низкой множественностью вторичных частиц и довольно малой потерей энергии снаряда. Трассировка лидера прекращается, когда его энергия падает ниже определенного порога — например, 2\% от первичной энергии. Это снижение энергии помогает обеспечить более сбалансированное распределение энергии между параллельными процессами. Хотя точный порог может варьироваться, цель состоит в том, чтобы снизить энергию лидера до уровня, который гарантирует генерацию достаточного количества вторичных частиц для начала эффективной параллельной трассировки.

Далее главный процесс разделяет стек вторичных частиц между параллельными процессами и начинает агрегировать черенковские фотоны, произведенные в подкаскадах, в многомерные массивы. CORSIKA генерирует и трассирует только те черенковские фотоны, которые производят фотоэлектроны в фотодатчике с единичной вероятностью. 

После того как все ведомые процессы завершают свои задачи, массивы сохраняются в результирующий двоичный файл. Такое распределение обязанностей между главным и ведомыми процессами показало свою высокую эффективность.


\subsection{Алгоритм распределения стека}

Результат распределения стека частиц между вторичными процессами имеет решающее значение для ускорения генерации событий. Когда трассировка лидера завершается, главный процесс вычисляет общее энергосодержание стека $E_{\text{tot}}$ и сортирует частицы в порядке возрастания энергии. Целевая энергия на один вторичный процесс рассчитывается как:
\[
e_p = \frac{E_{\text{tot}}}{n_{\text{slaves}}},
\]
где $n_{\text{slaves}}$ обозначает количество доступных вычислительных процессов.
Затем главный процесс последовательно формирует подстеки для каждого ведомого процесса. Частицы выбираются из отсортированного списка, начиная с наименее энергичных, и их индексы помещаются в массив подстека. Накапливается суммарная энергия $e_{\text{pp}}$ текущего подстека. Формирование подстека для текущего процесса завершается, как только $e_{\text{pp}}$ превышает $e_{\text{p}}$. Этот подход гарантирует, что энергия, назначенная каждому процессу, остается близкой к среднему значению, а результирующие подстеки содержат частицы со сравнимыми полными энергиями.
Однако на практике равномерное распределение не гарантируется: поскольку отдельные частицы неделимы, суммарная энергия в подстеке может либо превышать $e_{\text{pp}}$, либо значительно отставать от него. %Эта проблема особенно выражена, когда в стеке присутствуют гамма-кванты высокой энергии, поскольку их энергия иногда может составлять до 50\% от общей энергии. В таких случаях подстек, назначенный одному процессу, может накопить энергию, существенно превышающую целевое значение, оставляя недостаточную энергию для последующих процессов. В результате могут возникать ситуации, когда некоторые процессы не получают частиц вообще.

Такой прямой подход к распределению нагрузки явно требует оптимизации. Простым улучшением является расчет целевой энергии подстека как $E_{\text{tot}}/(n_{\text{slaves}} + 1)$, а не ${E_{\text{tot}}}/{n_{\text{slaves}}}$. Другим соображением является то, что порог снижения энергии частицы-лидера должен корректироваться в зависимости от количества процессов; например, если используется 10 ведомых процессов при запуске программы, то энергия частицы-лидера, составляющая 2\% от первичной энергии, теоретически может дать в одном из подстеков до 12\% от начальной энергии, хотя при идеальном разделении было бы 10\%. В таком случае максимальное увеличение энергии в одном из подстеков составляет до 20\% от расчетного значения. Но если мы хотим запустить программу, используя 100 ведомых процессов, то порог 2\% уже является слишком высоким: расчетная энергия на один процесс будет составлять 1\% от первичной, и процесс в таком случае может получить до 3\% первичной энергии, то есть максимальное количество энергии в подстеке увеличивается до 3 раз по сравнению с расчетной величиной, что влечет значительное увеличение работы программы. 

На данный момент программа запускается с использованием 10 ведомых процессов. Трассировка лидера завершается либо когда энергия лидера падает до 2\% от первичной энергии, либо когда в стеке появляется гамма-квант высокой энергии. Последнее условие обусловлено тем фактом, что электромагнитные каскады, инициированные гамма-квантами, требуют больших вычислительных затрат. Продолжение последовательной обработки в главном процессе в таких случаях становится неэффективным; более выгодно распределить стек между вторичными процессами и инициировать фазу параллельных вычислений. Порог энергии в 2\% был выбран эмпирически исходя из следующих соображений: снижение энергии до такой величины обходится небольшими вычислительными затратами, а погрешность размера стека до 20\% в б\'oльшую сторону не сильно ухудшает желаемое ускорение программы.

%для улучшения разделения энергии между вычислительными процессами, при этом минимизируя время, затрачиваемое на последовательную трассировку, и обеспечивая более сбалансированное распределение общей энергии.%
Будущие улучшения будут сосредоточены на оптимизации алгоритма разделения стека для достижения более равномерного распределения частиц между ведомыми процессами.


%\vspace{2cm}
\section{Результаты тестирования}
\subsection{Конфигурация тестовой платформы и экспериментальные параметры}
\label{subsec:test_platform}

Для отладки параллельного алгоритма и проведения предварительных тестов использовался локальный сервер со следующей конфигурацией:
один 16-ядерный процессор AMD Ryzen 9 5950X, базовая частота 3.4~ГГц;
материнская плата ASUS TUF GAMING X570-PRO, сокет AM4, чипсет AMD X570, поддерживающий до 32 потоков выполнения;
оперативная память 128~ГБ DDR4 ECC (4~модуля по 32~ГБ каждый, частота 3200~МГц);
операционная система Linux Fedora 38.

Конфигурация сервера обеспечивает высокую производительность одного потока и достаточный объем памяти для одновременной обработки нескольких событий ШАЛ без использования файла подкачки на диске. Генерация одного события потребляет примерно 7~ГБ памяти для главного процесса и около 450~МБ для каждого вторичного субпроцесса. Следовательно, с учетом системных резервов, сервер может одновременно выполнять до 15 событий в последовательной версии программы или два события в параллельной версии с использованием 10 дополнительных процессов.

Этот сервер использовался для разработки и тестирования параллельной версии программы, а также для сравнения ее результатов с результатами оригинальной версии. Сравнение проводилось на основе моделирования наборов тестовых событий со следующими параметрами:

\begin{itemize}
\item тип первичной частицы: протон и ядро железа ($^{56}$Fe);
\item энергии первичных частиц: $10^{15}$, $10^{16}$ и $10^{17}$~эВ;
\item зенитный угол: $\theta = 0^\circ$;
\item азимутальный угол: равномерное распределение в диапазоне 0--360$^\circ$;
\item модель адронных взаимодействий: QGSJET-II~\cite{qgsIIa, qgsIIb};
\item модель атмосферы: стандартная атмосфера США (номер~1 в пакете CORSIKA).
\end{itemize}

Для обеих версий программы было смоделировано по 50 событий для каждого набора параметров. Для тестирования параллельной версии количеством рабочих процессов было выбрано $n_{\text{slaves}}$=10.


\subsection{Результаты тестирования: производительность}
\label{subsec:experimental_results}

Для оценки производительности мы измерили время выполнения  программы на тестовом сервере.
%(см. рис.~\ref{fig:mean_time_server}). 
Для событий, инициированных первичными протонами с энергией $10^{17}$~эВ, среднее время обработки одного события снизилось с 21~часа в последовательной версии до 8.5~часов в параллельной версии. Полученное ускорение $S = T_{\mathrm{seq}} / T_{\mathrm{par}}$ для этого набора параметров равно 2.4. Значения ускорения для других типов первичных частиц и энергий показаны на рис.~\ref{fig:mean_time_server_relative} и варьируются от 1.95 до 3.26 в зависимости от энергии и ядерного состава.

%%\begin{figure}[ht]
%%\includegraphics[width=10cm]{figs/mean_time_server.eps}
%%\caption{Среднее время расчета одного события в двух версиях программы на локальном сервере.}\label{fig:mean_time_server}
%%\end{figure}

\begin{figure}[htb]
\includegraphics[width=8cm]{figs/mean_time_server_relative.eps}
\caption{Среднее ускорение расчета одного события для разных энергий при одинаковом числе вторичных процессов $p=10$}\label{fig:mean_time_server_relative}
\end{figure}

Зависимость фактора ускорения $S(p)=T_{seq} / T_{par} (p)$ от числа вычислительных процессов $p$ для параллельной версии кода приведено на рис.~\ref{fig:mean_speedup_multi_server}. Результаты получены при моделировании ШАЛ, инициированных ядрами железа с энергией $10^{17}$ эВ, на вычислительном сервере. Пунктирная линия соединяет экспериментальные точки для наглядности. Как видно из графика, масштабирование кода близко к линейному в диапазоне от 2 до 13 процессов, после чего наблюдается насыщение фактора ускорения вследствие увеличения накладных расходов на межпроцессорное взаимодействие и синхронизацию данных. При $n=1$ фактор ускорения меньше единицы, что обусловлено дополнительными затратами времени на инициализацию параллельной среды и распределение данных между процессами по сравнению с непараллельной версией.

\begin{figure}[htb]
\includegraphics[width=8cm]{figs/mean_speedup_multi_server.eps}
\caption{Средний фактор ускорения расчета одного события $S(p)$ в зависимости от  числа вычислительных процессов}\label{fig:mean_speedup_multi_server}
\end{figure}

Эффективностью распараллеливания $E(p)$ называется безразмерная величина,  показывающая, насколько рационально используются вычислительные ресурсы (процессорные ядра или узлы кластера) при решении задачи параллельным алгоритмом. Она определяется как отношение фактора ускорения $S(p)$ к числу задействованных вычислительных процессов $p$. Как следует из рис.~\ref{fig:mean_effectiv_multi_server}, для нашей версии программы максимальная эффективность $E$, достигающая $\sim$33\%, наблюдается при использовании 6 процессов. В диапазоне от 3 до 13 процессов эффективность остается относительно стабильной на уровне $\sim$30\%. Дальнейшее увеличение числа процессов приводит к постепенному снижению эффективности из-за возрастающих накладных расходов на межпроцессное взаимодействие и синхронизацию. Относительно низкое абсолютное значение эффективности (около 0.3) указывает, вероятнее всего, на наличие узких мест при обмене данными. %, что согласуется с фундаментальными ограничениями, описываемыми законом Амдала.


\begin{figure}[htb]
\includegraphics[width=8cm]{figs/mean_effectiv_multi_server.eps}
\caption{Зависимость эффективности распараллеливания $E(p)=S(p)/p$ от числа вычислительных процессов}\label{fig:mean_effectiv_multi_server}
\end{figure}


\subsection{Результаты тестирования: физическая валидация}
\label{subsec:experimental_results}

Диаграмма размаха числа взаимодействий частицы-лидера до достижения порога энергии в 2\% от первичной энергии представлена на рис.~\ref{fig:interactions}. Ядра железа демонстрируют немного большее среднее число взаимодействий по сравнению с протонами той же энергии (например, 11.3 против 7.1 при 10 ПэВ), это число  практически не зависит от энергии. Разброс значений (мин--макс) для протонов несколько превышает таковой для ядер железа. Такая изменчивость обусловлена флуктуациями глубины первого взаимодействия и стохастическим характером развития ядерного каскада, что более выражено для легких первичных частиц. 


\begin{figure}[bht]
%\setcaptionmargin{5mm}
%\onelinecaptionsfalse % if the caption is multiline
%\onelinecaptionstrue  % if the caption is one-line
\includegraphics[width=10cm]{figs/interaction_boxplot.eps}
%\captionstyle{normal}
\caption{Диаграмма размаха числа взаимодействий лидирующей частицы до достижения порога энергии в 2\%.}\label{fig:interactions}
\end{figure}


Валидация характеристик ШАЛ по двум версиям программы проведена на сервере для программы с $p=10$. Для проверки физической корректности параллельной реализации мы сравнили пространственные распределения черенковских фотонов на уровне наблюдения. Были проанализированы функции пространственного распределения (ФПР), усредненные по статистике событий и полученные для идентичных наборов параметров первичных частиц. На рисунке~\ref{fig:mean_ldf_proton_server} представлены усредненные поперечные распределения фотонов для первичных протонов трех энергий ($10^{15}$, $10^{16}$ и $10^{17}$~эВ), полученные для двух версий программного кода. Различия между последовательной и параллельной версиями незначительны и не превышают уровня статистических флуктуаций, ожидаемых для конечных размеров выборок. Систематического смещения не наблюдается.

\begin{figure}[htb]
%\setcaptionmargin{5mm}
%\onelinecaptionsfalse % if the caption is multiline
%\onelinecaptionstrue  % if the caption is one-line
\includegraphics[width=8cm]{figs/mean_ldf_proton_server.eps}
\includegraphics[width=8cm]{figs/mean_ldf_0014_server_relative.eps}
%\captionstyle{normal}
\caption{Сравнение средних функций пространственного распределения черенковских фотонов от первичного протона в двух версиях программы. Слева --- средние ФПР черенковского света, справа --- отношение ФПР, полученных с помощью параллельной версии, к оригинальной ФПР.} \label{fig:mean_ldf_proton_server}
\end{figure}

Оценка массы первичной частицы в эксперименте СФЕРА строится на отношении в ливне количества фотонов на разных расстояниях от оси ШАЛ. Поэтому важно, чтобы отношение ФПР различных ядер не искажалось при изменении способа расчета. На рис.~\ref{fig:mean_ldf_server_relative_to_Fe} видно, что отношение ФПР от протона к ФПР от ядра железа сохраняется в пределах статистических погрешностей при переходе к параллельной версии программы. 

\begin{figure}[htb]
\includegraphics[width=8cm]{figs/mean_ldf_nonparallel_server_relative_to_poster.eps}
\includegraphics[width=8cm]{figs/mean_ldf_parallel_server_relative_to_poster.eps}
\caption{Отношение ФПР черенковского света от протона к ФПР от ядра железа в двух версиях программы, слева --- исходной, справа --- параллельной.} \label{fig:mean_ldf_server_relative_to_Fe}
\end{figure}



Мы также сравнили среднее полное количество черенковских фотонов на одно событие. Для ядер железа относительное расхождение оставалось в пределах 2--6\% во всем диапазоне энергий $10^{15}-10^{17}$ эВ, для протонов расхождение было немного ниже, от 1\% до 4\%. Полученное согласие на уровне 1–6\% между последовательной и параллельной версиями кода свидетельствует о корректности реализованного алгоритма распараллеливания. Указанные расхождения находятся в пределах статистических флуктуаций, характерных для метода Монте-Карло, и не превышают систематических неопределенностей, свойственных самому коду CORSIKA.

%Эта разница объясняется двумя факторами. Во-первых, ливни, инициированные протонами в рассматриваемом диапазоне энергий, демонстрируют более высокую внутреннюю изменчивость: относительная дисперсия выхода фотонов значительна из-за флуктуаций глубины первого взаимодействия и развития электромагнитной каскадной компоненты. Во-вторых, разница в размерах статистических выборок между двумя версиями приводит к флуктуациям выборки в оценках средних значений. Для проверки отсутствия систематических ошибок, вносимых алгоритмом параллелизации, планируется дополнительный эксперимент с идентичными статистическими объемами.


Полученные результаты подтверждают, что примененная стратегия параллелизации сохраняет статистические свойства сгенерированных ансамблей и не вносит систематических искажений в физически значимые распределения.


\section{Заключение}

Мы разработали параллельную версию кода CORSIKA с опцией черенковского света для преодоления вычислительных ограничений, возникающих при моделировании широких атмосферных ливней при энергиях свыше $7\cdot10^{16}$~эВ на суперкомпьютере «Ломоносов-2». Реализация использует двухэтапный алгоритм: последовательная трассировка первичной частицы и ее наиболее энергичного потомка (лидера) до тех пор, пока его энергия не упадет ниже 2\% от первичной энергии, за которой следует параллельная обработка результирующего стека частиц, распределенного между вычислительными процессами.

Модификация кода CORSIKA, которую мы когда-то задумали для решения нашей технической проблемы с «Ломоносовым-2», теперь полностью работоспособна. Тесты производительности демонстрируют коэффициент ускорения $S = 1.95$--$3.26$ в зависимости от типа и энергии первичной частицы, сокращая время обработки событий от протонов с энергией $10^{17}$~эВ с 21~часа до 8.5~часов на тестовой серверной платформе. Достигнутое ускорение не является рекордным, но соответствует нашим ожиданиям.

Физическая валидация подтверждает статистическую эквивалентность последовательной и параллельной версий: функции пространственного распределения черенковских фотонов и средний выход фотонов согласуются в пределах 1--4\% для протонов и 1--6\% для ядер железа, что согласуется с ожидаемыми флуктуациями выборки.

Таким образом, разработанная параллельная версия CORSIKA воспроизводит результаты оригинального кода с точностью, достаточной для задач моделирования черенковского излучения ШАЛ в проекте СФЕРА-3, где типичные экспериментальные неопределенности значительно превышают наблюдаемые расхождения.

Основным ограничением текущей реализации является неравномерное распределение частиц между вычислительными процессами из-за природы взаимодействий частиц высоких энергий, что иногда приводит к простаиванию некоторых процессов. Будущая работа будет сосредоточена на оптимизации алгоритма разделения стека и исследовании ускорения на GPU в наших вычислениях, по крайней мере, для некоторых фрагментов кода. Несмотря на эти ограничения, разработанный код уже позволяет эффективно генерировать большие базы данных событий ШАЛ для проекта СПФЕРА-3 и может быть полезен для других экспериментов с космическими лучами, требующих моделирования черенковского света.

%\section{Благодарности}
\section*{Благодарности}

%\begin{acknowledgments}
%Исследование проводилось в рамках государственного задания Московского государственного университета имени М.В. Ломоносова, 
%The study was conducted under the state assignment of Lomonosov Moscow State University

Работа поддержана Российским научным фондом (грант № 23-72-00006). %https://rscf.ru/project/23-72-00006/), 

Работа В.А. Иванова поддержана Фондом поддержки теоретической физики и математики «БАЗИС» (грант № 24-2-10-53-1).

Исследование выполнено с использованием оборудования Центра коллективного пользования суперкомпьютерными ресурсами Московского государственного университета имени М.В. Ломоносова~\cite{SC}.
%\end{acknowledgments}
	
%%  ===================================================================
\begin{thebibliography}{20}

\bibitem{KAMPERT2012660}
K.~Kampert and M.~Unger, Astropart. Phys. \textbf{35}, 660 (2012).
%K. Kampert and M. Unger / {Measurements of the cosmic ray composition with air shower experiments} // {Astroparticle Physics}. 2012. V.35. 10. P. 660-678.
\\ https://doi.org/10.1016/j.astropartphys.2012.02.004

        
\bibitem{Chudakov1972}
%\refitem{proceedings}
А.~Е.~Чудаков, Экспериментальные методы исследования космических лучей сверхвысоких энергий: Материалы Всесоюзного симпозиума,
Якутск, 19–23 июня 1972 (Якутск. фил. Сиб. отд. АНСССР, 1974), с. 69.
%A.~E.~Chudakov, \textquotedblleft Probable method for EAS registration using Cherenkov light reflected from snow surface\textquotedblright~[in russian], Proc. All-Union Symp. 69--74 (1972).

\bibitem{Chernov2015}
Д.В. Чернов, Р.А Антонов., Т.В. Аулова, Е.А. Бонвеч, и др. // ЭЧАЯ. \textbf{46} (1) 115 (2015). [D.V. Chernov, R.A. Antonov, T.V. Aulova, E.A. Bonvech, et. al. // Phys. Part. Nucl. \textbf{46} (1) 60 (2015). http://doi.org/10.1134/S1063779615010025 ]


\bibitem{2020_SPHERE_JINST}
D.~Chernov, E.~Bonvech, M.~Finger et al., JINST \textbf{15}, C09061 (2020).
%\textquotedblleft Investigation of the energy spectrum and chemical composition of primary cosmic rays in 1--100~PeV energy range with a UAV-borne detector\textquotedblright, JINST, \textbf{15}, C09061 (2020).
\\ http://doi.org/10.1088/1748-0221/15/09/C09061


\bibitem{Chernov2024}
D.~V.~Chernov et al., Phys. At. Nucl. \textbf{87} (S2), S319 (2024).
%\textit{Chernov D.V. et al.} // Phys. At. Nucl. 2024. V. 87 (S2). P. S319. \\
\\ https://doi.org/10.1134/S1063778824700959


\bibitem{mass}
%\refitem{article}
V.~I.~Galkin, C.~G.~Azra, E.~A.~Bonvech et al., 
%\textquotedblleft SPHERE-3: Tackling the Problem of Primary Cosmic Ray Mass Composition with a New Approach\textquotedblright, Moscow University Physics Bulletin, \textbf{79}(S1), P. S379--S386 (2024).
Moscow Univ. Phys. Bull. \textbf{79} (S1), S379 (2024).
http://doi.org/10.3103/S002713492470111X


\bibitem{CORSIKA-7}
%\refitem{report}
D.~Heck, J.~Knapp, J.~N.~Capdevielle et al., 
%\textquotedblleft CORSIKA: a Monte Carlo code to simulate extensive air showers\textquotedblright, 
Forschungszentrum Karlsruhe Report \textbf{FZKA 6019} (1998).
http://bibliothek.fzk.de/zb/berichte/FZKA6019.pdf


\bibitem{SC}
V.~V.~Voevodin, A.~S.~Antonov, D.~A.~Nikitenko, P.~A.~Shvets et al., Supercomput. Front. Innov. \textbf{6} (2), 4 (2019).
https://doi.org/10.14529/jsfi190201


\bibitem{2025NuclPhys}
E.~A.~Bonvech, O.~V.~Cherkesova, D.~V.~Chernov, E.~L.~Entina et al., Phys. At. Nucl. \textbf{88} (12), 2522 (2025).
https://doi.org/10.1134/S1063778825100023

\bibitem{qgsIIa}
S.~S.~Ostapchenko, Phys. Rev. D \textbf{83}, 014018 (2011).
%\textquotedblleft Monte Carlo treatment of hadronic interactions in enhanced Pomeron scheme: QGSJET-II model \textquotedblright, {\it Phys. Rev.} {\bf D83} (2011) 014018.
\\ https://doi.org/10.1103/PhysRevD.83.014018


\bibitem{qgsIIb} 
S.~S.~Ostapchenko, Phys. Rev. D \textbf{89}, 074009 (2014).\\
%S.S. Ostapchenko, \textquotedblleft LHC data on inelastic diffraction and uncertainties in the predictions for longitudinal extensive air shower development \textquotedblright, {\it Phys. Rev.} {\bf D89} (2014) 074009.
https://doi.org/10.1103/PhysRevD.89.074009

\end{thebibliography}
	
	
	
	% ======================================
	% Перевод шапки статьи на английский
	% ======================================
	\newpage
	\begin{center}
		\textbf{EAS Cherenkov light efficient simulation: CORSIKA modification for the SPHERE-3 project}
		
		\medskip\small
		
		\textbf{E.\,A. Bonvech$^{1,a}$, M.\,D. Ziva$^{1,2}$, O.\,V. Cherkesova$^{1,3}$, D. V. Chernov$^{1}$, V.\,I. Galkin$^{1,4}$, V.\,A. Ivanov$^{1,4}$, T.\,A. Kolodkin$^{1,4}$, N.\,O. Ovcharenko$^{1,4}$, D.\,A. Podgrudkov$^{1,4}$, T.\,M. Roganova$^1$}
		
		\smallskip
		
		\textit{$^1$ Skobeltsyn Institute for Nuclear Physics, Lomonosov Moscow State University}\\
		\textit{$^2$ Faculty of Computational Mathematics and Cybernetics, Lomonosov Moscow 
        \textit{$^3$ Department of Space Research, Lomonosov Moscow State University}\\
		State University}\\
        \textit{$^4$ Faculty of Physics, Lomonosov Moscow State University}\\
        
		\textit {E-mail: $^a$ bonvech@yandex.ru}
		
		\medskip
		Abstract
		\smallskip
		
		\textit{Keywords}: extensive air showers, detector, Cherenkov light, simulation\\
		
		
	\end{center}
	
	
\end{document}

%%%%%%%%%%%%%%%%%%%%%%%%%%% Character code reference %%%%%%%%%%%%%%%%%%%%%%%%%%
%                                                                             %
%     Upper case russian letters (CP866): АБВГДЕЁЖЗИЙКЛМНОПРСТУФХЦЧШЩЪЫЬЭЮЯ   %                                                                       %
%     Lower case russian letters (CP866): абвгдеёжзийклмнопрстуфхцчшщъыьэюя   %
%                     Upper case letters: ABCDEFGHIJKLMNOPQRSTUVWXYZ          %
%                     Lower case letters: abcdefghijklmnopqrstuvwxyz          %
%                                   Digits: 0123456789                        %
% Square, curly, angle braces, parentheses: [] {} <> ()                       %
%                Backslash, slash, solidus: \ / |                             %
%       Period, interrogative, exclamation: . ? !                             %
%                 Comma, colon, semi-colon: , : ;                             %
%          Underscore, hyphen, equals sign: _ - =                             %
%             Quotes (left, right, double): ` ' "                             %
%     Commercial-at, hash, dollar, percent: @ # $ %                           %
%  Ampersand, asterisk, plus, caret, tilde: & * +   ^                         %
%                                                                             %
%%%%%%%%%%%%%%%%%%%%%%%%%%%%%%%%%%%%%%%%%%%%%%%%%%%%%%%%%%%%%%%%%%%%%%%%%%%%%%%
%
% Пожалуйста, включайте в свой файл выше приведенные строки
% c "Character code reference", чтобы была возможность
% проконтролировать правильность пересылки Вашего файла по e-mail .
%
% Конец файла paper.tex
%